\scanpage{42}
\begin{center}-- 25 -- \hfill n° 371\end{center}
lois externes sur $M$. Soit $\Omega$ l'ensemble \uline{somme} (\uline{Ens.}, chap. II, § 4, n° 8) des domaines d'opérateurs des diverses lois externes de $M$, chacun de ces domaines d'opérateurs étant identifié canoniquement à la partie de $\Omega$ qui lui correspond. Soit $L(\Omega)$ le \uline{monoïde libre} (chap. I, 3\textsuperscript{e} éd., § 1, n° 3) déduit de $\Omega$; nous définirons une loi de composition externe $(\alpha,x)\to\alpha x$ sur $M$, ayant $L(\Omega)$ comme domaine d'opérateurs, par récurrence sur la longueur du \uline{mot} $\alpha$ dans $L(\Omega)$; si $\alpha$ est de longueur $0$, c'est le mot vide $e$, et nous posons $ex=x$ pour tout $x\in M$. Si maintenant $\alpha$ est de longueur $n\geq1$, on peut l'écrire d'une seule manière sous la forme $\alpha=\beta\gamma$, où $\gamma$ est de longueur $n-1$ et $\beta$ de longueur $1$; $\beta$ appartient par suite à un et un seul des domaines d'opérateurs des lois externes données sur $M$ et $\beta y$ est donc défini pour tout $y\in M$; on pose alors $\alpha x=\beta(\gamma x)$. Pour deux mots quelconques $\alpha$, $\beta$ dans $L(\Omega)$, on vérifie que l'on a $\alpha(\beta x)=(\alpha\beta)x$ pour tout $x\in M$, par récurrence sur la longueur de $\alpha$.

Soit maintenant $A$ l'algèbre du monoïde $L(\Omega)$ relative à l'anneau $Z$ des entiers rationnels; on définit une loi externe $(\alpha,x)\to\alpha x$ sur $M$, ayant $A$ pour anneau d'opérateurs, en posant, pour tout $\alpha=\sum_{\lambda\in L(\Omega)}n_\lambda.\lambda$ (avec $n_\lambda\in Z$ et $n_\lambda=0$ sauf pour un nombre fini de valeurs de $\lambda$),
\[
\alpha x=\sum_{\lambda\in L(\Omega)}n_\lambda(\lambda x).
\]
On vérifie sans peine que cette loi externe satisfait aux axiomes $(\unsure{M_{\mathrm I}|N_{\mathrm I}})$, $(M_{\mathrm{II}})$, $(M_{\mathrm{III}})$ et $(M_{\mathrm{IV}})$ du chap. II, 3\textsuperscript{e} éd., § 1, n° 1; elle définit donc sur $M$ une structure de $A$-module à gauche, et il est immédiat que cette structure satisfait aux conditions 1) et 2) ci-dessus.

La méthode précédente s'applique à des groupes commutatifs à opérateurs quelconques; mais si on ne considère que des groupes commutatifs dont les opérateurs satisfont à certaines conditions supplémentaires, on peut souvent satisfaire aux conditions 1) et 2) en associant d'une autre manière une structure de module à la structure donnée de groupe à opérateurs sur $M$.

Par exemple, supposons qu'il n'y ait sur $M$ qu'une seule loi externe, dont le domaine d'opérateurs soit \uline{une partie} $S$ \uline{d'un anneau} $B$; on sait
\mspage{41}
